% Expediente para la reunión con el abogado: qué hace el sistema con los
% datos, en palabras de no técnicos, y las preguntas que hay que contestar.
% Sale de legal/LEEME.md, /privacidad y el contrato de encargo: si cambian,
% cambia esto.
\documentclass[10.5pt,a4paper]{article}
% Preámbulo común de los documentos legales (legal/*.tex).
% Se compila cada uno con:   cd legal && latexmk -pdf <archivo>.tex
% La paleta es la del producto y la del manual: si cambia una, cambian las dos.
\usepackage[T1]{fontenc}
\usepackage[utf8]{inputenc}
\usepackage[spanish,es-noshorthands,es-tabla]{babel}
\usepackage[scaled=0.92]{helvet}
\renewcommand{\familydefault}{\sfdefault}
\usepackage{amssymb}
% Los que usan los textos y pdflatex no trae
\DeclareUnicodeCharacter{2713}{\checkmark}
\DeclareUnicodeCharacter{22EF}{\ensuremath{\cdots}}
\DeclareUnicodeCharacter{2192}{\ensuremath{\rightarrow}}
\usepackage[a4paper,top=2cm,bottom=2.2cm,left=2.2cm,right=2.2cm,headheight=14pt]{geometry}
\usepackage[table]{xcolor}
\usepackage{microtype}
\usepackage{enumitem}
\usepackage{booktabs}
\usepackage{tabularx}
\usepackage{array}
\usepackage{longtable}
\usepackage{fancyhdr}
\usepackage{titlesec}
\usepackage{lastpage}
\usepackage[most]{tcolorbox}
\usepackage[hidelinks]{hyperref}

\definecolor{tinta}{HTML}{0B2239}
\definecolor{tinta2}{HTML}{2A4864}
\definecolor{tenue}{HTML}{496A86}
\definecolor{marca}{HTML}{1668C4}
\definecolor{ambar}{HTML}{A67300}
\definecolor{rojo}{HTML}{C2001D}
\definecolor{fondo}{HTML}{F4F7FB}
\definecolor{linea}{HTML}{C3D4E4}
\definecolor{ambarfondo}{HTML}{FDF6E4}
\definecolor{ambarborde}{HTML}{E4C88A}
\definecolor{azulfondo}{HTML}{E7F0FB}
\definecolor{azulborde}{HTML}{C9DDF4}

\color{tinta}
\setlength{\parindent}{0pt}
\setlength{\parskip}{5pt plus 2pt}
\linespread{1.08}
\setlist{itemsep=2pt,topsep=3pt}

\titleformat{\section}{\color{tinta}\large\bfseries}{}{0em}{}
\titleformat{\subsection}{\color{tinta2}\normalsize\bfseries}{}{0em}{}
\titlespacing*{\section}{0pt}{12pt}{4pt}
\titlespacing*{\subsection}{0pt}{8pt}{2pt}

\renewcommand{\arraystretch}{1.25}
\arrayrulecolor{linea}

% Lo que hay que completar a mano: una línea para escribir
\newcommand{\llenar}[1][6cm]{\rule[-2pt]{#1}{0.5pt}}
% Lo que tiene que decidir o completar el abogado o las partes
\newcommand{\completar}[1]{{\color{ambar}[#1]}}
% Una casilla para marcar
\newcommand{\casilla}{\fbox{\rule{0pt}{7pt}\rule{7pt}{0pt}}}

\newtcolorbox{nota}[1][]{colback=azulfondo,colframe=azulborde,boxrule=0.6pt,arc=3pt,
  left=8pt,right=8pt,top=5pt,bottom=5pt,fontupper=\small\color{tinta2},#1}
\newtcolorbox{borrador}{colback=ambarfondo,colframe=ambarborde,boxrule=0.6pt,arc=3pt,
  left=8pt,right=8pt,top=5pt,bottom=5pt,fontupper=\small}

% Pie: qué documento es, su versión y la página
\newcommand{\documento}[2]{%
  \pagestyle{fancy}\fancyhf{}%
  \renewcommand{\headrulewidth}{0pt}%
  \fancyfoot[L]{\scriptsize\color{tenue}#1 · #2}%
  \fancyfoot[R]{\scriptsize\color{tenue}Página \thepage\ de \pageref*{LastPage}}%
}

% Firma: una línea y lo que va debajo
\newcommand{\firma}[1]{\begin{minipage}[t]{0.44\textwidth}\vspace{1.6cm}\rule{\linewidth}{0.5pt}\\{\small\color{tinta2}#1}\end{minipage}}

\begin{document}
\documento{Protección de datos personales — expediente para el abogado}{11/11/2026}

{\LARGE\bfseries Protección de datos personales}\par
{\large\color{tinta2}Expediente para la revisión legal del sistema de control de flota «Control de ruta»}\par
{\small\color{tenue}Preparado el 11/11/2026 · Ley N.º 29733 y su Reglamento (D.S. N.º 016-2024-JUS)}

\begin{nota}
\textbf{Para qué es este documento.} Resume, sin tecnicismos, qué datos
personales trata el sistema, de quién, para qué, cuánto tiempo y dónde, y
termina con las \textbf{preguntas concretas} que necesitamos que usted
conteste (sección 9). Se adjuntan los borradores a revisar: el contrato de
encargo, el formulario de consentimiento, la política de privacidad, el
registro de consentimientos y la solicitud de derechos.
\end{nota}

\section{1. Qué es el sistema}
Un servicio que contratan cooperativas de transporte urbano (combis) de
Juliaca y Puno para controlar su flota. Tiene tres partes:
\begin{itemize}[leftmargin=14pt]
\item \textbf{La app del conductor y del cobrador} («Control de ruta», Android):
  mientras el conductor declara que está en ruta, manda su ubicación para
  medir la separación con la combi de adelante y la de atrás; tiene chat con
  la central y un botón de emergencia (SOS).
\item \textbf{El panel de la central} (Despacho y gerencia de cada cooperativa):
  ve las combis en el mapa, atiende los SOS, chatea, da de alta a su gente y
  saca informes de vueltas y horas trabajadas.
\item \textbf{El panel del operador} (nosotros): da de alta las cooperativas,
  carga los recorridos y mantiene el servicio.
\end{itemize}
Las cuentas las crea cada cooperativa para su propia gente: no hay registro
abierto ni usuarios anónimos. Cada cooperativa ve sólo sus datos.

\section{2. Quién es quién}
\begin{tabularx}{\linewidth}{>{\raggedright}p{4.2cm}>{\raggedright}p{3.6cm}X}
\toprule
\textbf{Rol en la ley} & \textbf{Quién} & \textbf{Por qué}\\\midrule
Titular de los datos & Cada conductor y cobrador (y el personal de la central) & Son sus datos: su ubicación, su nombre, sus horas\\
Responsable del tratamiento (titular del banco de datos) & Cada cooperativa & Ella decide controlar su flota y para qué; inscribe el banco, junta los consentimientos y responde los reclamos\\
Encargado del tratamiento & El operador del sistema & Guarda y procesa los datos por cuenta de la cooperativa y según sus instrucciones; lo obliga el contrato de encargo\\
Sub-encargados & Railway (servidor), el almacenamiento de las copias de respaldo, Geoapify (mapas) & Proveedores sin los cuales el servicio no funciona\\
\bottomrule
\end{tabularx}

\textbf{Un dato que cambia las respuestas:} muchos conductores no son
empleados de la cooperativa, sino socios o dueños de su combi.

\section{3. Qué datos, y cuándo}
\begin{itemize}[leftmargin=14pt]
\item \textbf{Ubicación durante el turno}, continua, \textbf{también con la
  pantalla apagada o la app en segundo plano} (Android lo exige declarado y
  Google Play lo revisa). Fuera del turno no se manda.
\item \textbf{Antes de empezar el turno}, con la app abierta, el teléfono
  \emph{mira} la ubicación sin mandarla: si ve que el conductor ya está
  manejando y no lo declaró, le avisa y, si no contesta en un minuto, lo pone
  en ruta solo (la cooperativa ve que fue automático).
\item \textbf{Lugares puntuales del turno}, con su hora: dónde se declaró en
  ruta, ausente o fuera ---\textbf{incluido dónde terminó el turno, que suele
  ser el domicilio del conductor}---, dónde se cortó la señal, dónde estuvo
  parado y dónde salió del recorrido.
\item \textbf{Marcas automáticas}: el sistema anota por su cuenta señales
  como «GPS sospechoso» (un GPS que parece simulado), «entrada tardía» (se
  sumó a la ruta por la mitad), «salió solo» o «fuera de ruta». \textbf{No
  producen ninguna consecuencia automática}: son avisos para que una persona
  de la cooperativa los mire. El contrato obliga a la cooperativa a no
  sancionar sólo por una marca y a explicarla si se la piden.
\item \textbf{Identidad}: nombre, alias, usuario, rol, combi.
\item \textbf{Chat}: textos, notas de voz y fotos que la persona decide
  mandar; las emergencias (SOS) con su ubicación.
\item \textbf{Horas trabajadas} (turnos) y \textbf{vueltas} cumplidas: con
  esto la cooperativa liquida el trabajo.
\item \textbf{Estado del teléfono} en ruta (batería, ahorro de batería,
  versión de la app): sólo el último valor, no se guarda.
\end{itemize}
No hay datos sensibles en el sentido de la ley (salud, origen, creencias…),
salvo lo que alguien escriba en el chat.

\section{4. Cuánto tiempo se guardan}
El propio sistema los borra solo al cumplirse el plazo:

\begin{tabularx}{\linewidth}{>{\raggedright}p{8.4cm}X}
\toprule
Chat (texto, fotos, notas de voz) & 30 días\\
Mensajes que la central guardó como evidencia de un hecho & 365 días\\
Vueltas, tramos y desvíos & 120 días\\
Lugares puntuales del turno y marcas automáticas & 120 días\\
Turnos (horas) y emergencias (SOS) & 365 días\\
Registro de actividad & accesos 90 días; operación 365 días; altas, bajas y claves 2 años\\
Posición en vivo y estado del teléfono & sólo el último valor\\
Copias de respaldo cifradas & 30 días\\
\bottomrule
\end{tabularx}

\section{5. Dónde están}
En servidores \textbf{fuera del Perú}: Railway (servidor y base de datos,
\completar{región a confirmar}), el almacenamiento de las copias de respaldo
(\completar{proveedor y región}) ---que recibe las copias \textbf{ya cifradas},
sin poder leerlas--- y Geoapify (imágenes del mapa: recibe la dirección IP del
teléfono, no la identidad). Google Play distribuye la app.

\section{6. Seguridad}
Todo viaja cifrado (HTTPS). Las contraseñas y las sesiones se guardan como
huellas irreversibles: ni el operador ni quien robe una copia de la base
puede verlas ni usarlas. Cada cooperativa está separada de las demás, y
dentro de ella cada rol ve lo que le toca. Las acciones administrativas
(altas, bajas, claves, correcciones de horas) quedan registradas con quién y
cuándo. Hay respaldos automáticos verificados y una copia fuera del servidor,
cifrada con AES-256.

\section{7. Los derechos, en la práctica}
\begin{itemize}[leftmargin=14pt]
\item \textbf{Acceso y portabilidad}: la central baja un archivo con todo lo de
  una persona (mensajes, fotos, turnos, vueltas, actividad), con una guía de
  qué es cada cosa.
\item \textbf{Rectificación}: nombre y alias desde el panel; las horas se
  corrigen con motivo, y queda registrado.
\item \textbf{Cancelación}: la baja borra la cuenta y los mensajes privados;
  \textbf{lo que es registro de trabajo} (turnos, vueltas, actividad) queda
  hasta su plazo.
\item \textbf{Revocación del consentimiento}: se anota en el sistema y en papel.
  El sistema no apaga la app solo: sin la ubicación no cumple su función, y
  qué hacer es decisión de la cooperativa.
\item \textbf{Revisión de una marca automática}: la central la ve con su
  detalle y hora, y puede explicarla.
\end{itemize}

\section{8. Si hay un incidente}
El contrato obliga al operador a avisar a la cooperativa \textbf{dentro de las
24 horas}, para que ella cumpla las \textbf{48 horas} ante la Autoridad y,
cuando corresponda, ante los titulares.

\section{9. Lo que necesitamos que conteste}
\begin{enumerate}[leftmargin=18pt,itemsep=4pt]
\item ¿Está bien el reparto de roles de la sección 2 (cooperativa
  responsable, operador encargado)?
\item La ubicación de un trabajador para controlar su trabajo, ¿necesita
  consentimiento o entra en la excepción de la relación laboral o
  contractual (art. 14)? ¿Cambia si el conductor es \textbf{socio o dueño de
  la combi} y no empleado? (Aunque no haga falta, preferimos tenerlo
  firmado.)
\item ¿Alcanza el formulario en papel, o conviene además que la app lo pida
  al entrar?
\item Los servidores fuera del Perú: ¿qué hay que declarar y qué garantías
  pedirle a cada proveedor (flujo transfronterizo)? ¿Se declara al inscribir
  el banco de datos o aparte?
\item ¿El operador tiene que inscribir un banco de datos propio (los
  contactos de cada cooperativa y de quien escribe a soporte)?
\item ¿Hay que designar un oficial de datos personales? ¿Cuenta como «gran
  volumen» la ubicación continua de cientos de conductores?
\item Las \textbf{marcas automáticas}: ¿entran en el derecho a un tratamiento
  objetivo y en lo que el Reglamento dice de las decisiones basadas sólo en
  tratamiento automatizado? ¿Alcanza con que no produzcan consecuencias solas
  y que el contrato prohíba sancionar sólo por ellas?
\item Los \textbf{lugares puntuales del fin de turno} (suele ser el domicilio),
  guardados 120 días: ¿es proporcional, o conviene acortarlo o guardar sólo
  la hora?
\item Al dar de baja a una persona, ¿está bien conservar el registro de trabajo
  (turnos, vueltas, actividad) hasta su plazo como excepción a la cancelación?
\item Si alguien retira el consentimiento, ¿qué tiene que hacer la cooperativa?
\item Los plazos para contestar cada derecho, para ponerlos en la solicitud.
\item Revisar y ajustar el contrato de encargo, el consentimiento (versión
  2026-11-11), la política de privacidad y la solicitud de derechos.
\end{enumerate}

\section{Documentos adjuntos}
\begin{itemize}[leftmargin=14pt,itemsep=0pt]
\item \texttt{contrato-encargo.pdf} — el contrato de encargo con cada cooperativa (borrador).
\item \texttt{consentimiento.pdf} — el formulario de consentimiento del conductor y del cobrador, versión 2026-11-11.
\item \texttt{politica-privacidad.pdf} — la política de privacidad publicada en la app, vigente desde el 11/11/2026.
\item \texttt{registro-consentimientos.pdf} — la hoja con que la cooperativa lleva las firmas y los retiros.
\item \texttt{solicitud-derechos.pdf} — el formulario para pedir acceso, rectificación, cancelación, oposición o portabilidad, o retirar el consentimiento.
\end{itemize}
\end{document}
